% --- Archon physics formalization source begin ---
% archon:physics
% archon:covers PhyXMiniProblems/problem_phyx_mini_0117.lean
% archon:source-report reports/phyx_mini/problem_phyx_mini_0117.source.json

\chapter{Physics problem phyx\_mini\_0117}
\label{ch:PhyXMiniProblems_problem_phyx_mini_0117}

\paragraph{Problem source.}
\#\# Physical scenario

Two radio antennas radiating in phase are located at points A and B, 200 m apart (figure). The radio waves have a frequency of 5.80 MHz. A radio receiver is moved out from point B along a line perpendicular to the line connecting A and B (line BC shown in figure).

\#\# Question

At what distances from B will there be destructive interference?

\#\# Answer choices

- A: A: 200m

- B: B: 250m

- C: C: 275m

- D: D: 225m

\#\# Figure caption (auxiliary; use the image as primary evidence)

The image contains a diagram with two vertical and one horizontal straight line intersecting. Points labeled \textbackslash{}( A \textbackslash{}) and \textbackslash{}( B \textbackslash{}) are placed along the vertical line, with point \textbackslash{}( B \textbackslash{}) below \textbackslash{}( A \textbackslash{}). There is a label indicating that the distance between \textbackslash{}( A \textbackslash{}) and \textbackslash{}( B \textbackslash{}) is \textbackslash{}( 200 \textbackslash{}, \textbackslash{}text\{m\} \textbackslash{}). Both points are marked by orange circles.

The vertical line, where \textbackslash{}( A \textbackslash{}) and \textbackslash{}( B \textbackslash{}) are located, intersects with a horizontal line extending to the right. This horizontal line has a point labeled \textbackslash{}( C \textbackslash{}) at the end. The labels \textbackslash{}( A \textbackslash{}), \textbackslash{}( B \textbackslash{}), and \textbackslash{}( C \textbackslash{}) denote different positions along the lines.

\#\# Dataset metadata

Wave Optics; Physical Model Grounding Reasoning, Spatial Relation Reasoning

\paragraph{Recorded answer/context.}
A: A: 200m

\paragraph{Figure/image path.}
/root/proposal\_for\_physic/science-mango/phyx\_mini\_run/phyx\_data/test\_image/117.png

\paragraph{Formalization target.}
create a compiling Lean file with sorry bodies at `PhyXMiniProblems/problem\_phyx\_mini\_0117.lean`. The Lean declarations must preserve the physical quantities, dimensions or dimensional roles, figure labels, governing-law hypotheses, and final relation expressed by this problem.
Use Mathlib/Physlib names found through LeanExplore where available. If a physics API is missing, introduce faithful local abstractions rather than scalar placeholder aliases.

\begin{theorem}[Physics formalization target]
\label{thm:physics:phyx_mini_0117:target}
\lean{PhyXMiniProblems.ProblemPhyXMini0117.problem_phyx_mini_0117}
\uses{def:physics:phyx-mini-0117:phyxminiproblems-problemphyxmini0117-twoantennainterferencesetup, def:physics:phyx-mini-0117:phyxminiproblems-problemphyxmini0117-radiatesinphaseatcommonfrequency, def:physics:phyx-mini-0117:phyxminiproblems-problemphyxmini0117-hasdepictedgeometry, def:physics:phyx-mini-0117:phyxminiproblems-problemphyxmini0117-hasstatedreadouts, def:physics:phyx-mini-0117:phyxminiproblems-problemphyxmini0117-hasphysicalparameters, def:physics:phyx-mini-0117:phyxminiproblems-problemphyxmini0117-satisfiesfreespacetwosourcelaws}
The assigned autoformalize agent should translate this physics problem into Lean declarations in the covered file, with theorem and lemma proof bodies written as `by sorry`.
\end{theorem}
\begin{proof}
This is an autoformalization task, not a proof task. Produce statements that can later be proved by the physics prover without weakening the physical model.
\end{proof}

% --- Archon declaration topology begin ---
\section{Declaration topology}
\begin{definition}[Length Quantity]
  \label{def:physics:phyx-mini-0117:phyxminiproblems-problemphyxmini0117-lengthquantity}
  \lean{PhyXMiniProblems.ProblemPhyXMini0117.LengthQuantity}
  A physical quantity carrying the dimension of length
\end{definition}
\begin{proof}
  This is the declared model definition of Length Quantity.
\end{proof}
\begin{definition}[Frequency Quantity]
  \label{def:physics:phyx-mini-0117:phyxminiproblems-problemphyxmini0117-frequencyquantity}
  \lean{PhyXMiniProblems.ProblemPhyXMini0117.FrequencyQuantity}
  A physical frequency, carrying the inverse-time dimension
\end{definition}
\begin{proof}
  This is the declared model definition of Frequency Quantity.
\end{proof}
\begin{definition}[Speed Quantity]
  \label{def:physics:phyx-mini-0117:phyxminiproblems-problemphyxmini0117-speedquantity}
  \lean{PhyXMiniProblems.ProblemPhyXMini0117.SpeedQuantity}
  A physical propagation speed
\end{definition}
\begin{proof}
  This is the declared model definition of Speed Quantity.
\end{proof}
\begin{definition}[meters Value]
  \label{def:physics:phyx-mini-0117:phyxminiproblems-problemphyxmini0117-metersvalue}
  \lean{PhyXMiniProblems.ProblemPhyXMini0117.metersValue}
  \uses{def:physics:phyx-mini-0117:phyxminiproblems-problemphyxmini0117-lengthquantity}
  The scalar readout of a physical length in metres
\end{definition}
\begin{proof}
  This is the declared model definition of meters Value.
\end{proof}
\begin{definition}[hertz Value]
  \label{def:physics:phyx-mini-0117:phyxminiproblems-problemphyxmini0117-hertzvalue}
  \lean{PhyXMiniProblems.ProblemPhyXMini0117.hertzValue}
  \uses{def:physics:phyx-mini-0117:phyxminiproblems-problemphyxmini0117-frequencyquantity}
  The scalar readout of a physical frequency in hertz
\end{definition}
\begin{proof}
  This is the declared model definition of hertz Value.
\end{proof}
\begin{definition}[megahertz Value]
  \label{def:physics:phyx-mini-0117:phyxminiproblems-problemphyxmini0117-megahertzvalue}
  \lean{PhyXMiniProblems.ProblemPhyXMini0117.megahertzValue}
  \uses{def:physics:phyx-mini-0117:phyxminiproblems-problemphyxmini0117-frequencyquantity}
  The scalar readout of a physical frequency in megahertz
\end{definition}
\begin{proof}
  This is the declared model definition of megahertz Value.
\end{proof}
\begin{definition}[meters Per Second Value]
  \label{def:physics:phyx-mini-0117:phyxminiproblems-problemphyxmini0117-meterspersecondvalue}
  \lean{PhyXMiniProblems.ProblemPhyXMini0117.metersPerSecondValue}
  \uses{def:physics:phyx-mini-0117:phyxminiproblems-problemphyxmini0117-speedquantity}
  The scalar readout of a physical speed in metres per second
\end{definition}
\begin{proof}
  This is the declared model definition of meters Per Second Value.
\end{proof}
\begin{definition}[Figure Point]
  \label{def:physics:phyx-mini-0117:phyxminiproblems-problemphyxmini0117-figurepoint}
  \lean{PhyXMiniProblems.ProblemPhyXMini0117.FigurePoint}
  The three point labels printed in the primary figure
\end{definition}
\begin{proof}
  This is the declared model definition of Figure Point.
\end{proof}
\begin{definition}[Antenna Label]
  \label{def:physics:phyx-mini-0117:phyxminiproblems-problemphyxmini0117-antennalabel}
  \lean{PhyXMiniProblems.ProblemPhyXMini0117.AntennaLabel}
  The two radiating antennas
\end{definition}
\begin{proof}
  This is the declared model definition of Antenna Label.
\end{proof}
\begin{definition}[figure Point]
  \label{def:physics:phyx-mini-0117:phyxminiproblems-problemphyxmini0117-antennalabel-figurepoint}
  \lean{PhyXMiniProblems.ProblemPhyXMini0117.AntennaLabel.figurePoint}
  \uses{def:physics:phyx-mini-0117:phyxminiproblems-problemphyxmini0117-figurepoint, def:physics:phyx-mini-0117:phyxminiproblems-problemphyxmini0117-antennalabel}
  The figure point occupied by each antenna
\end{definition}
\begin{proof}
  This is the declared model definition of figure Point.
\end{proof}
\begin{definition}[Answer Choice]
  \label{def:physics:phyx-mini-0117:phyxminiproblems-problemphyxmini0117-answerchoice}
  \lean{PhyXMiniProblems.ProblemPhyXMini0117.AnswerChoice}
  Labels of the four answers printed with the problem
\end{definition}
\begin{proof}
  This is the declared model definition of Answer Choice.
\end{proof}
\begin{definition}[displayed Distance Meters]
  \label{def:physics:phyx-mini-0117:phyxminiproblems-problemphyxmini0117-displayeddistancemeters}
  \lean{PhyXMiniProblems.ProblemPhyXMini0117.displayedDistanceMeters}
  \uses{def:physics:phyx-mini-0117:phyxminiproblems-problemphyxmini0117-answerchoice}
  The metre readout printed beside each answer label
\end{definition}
\begin{proof}
  This is the declared model definition of displayed Distance Meters.
\end{proof}
\begin{definition}[Radio Antenna]
  \label{def:physics:phyx-mini-0117:phyxminiproblems-problemphyxmini0117-radioantenna}
  \lean{PhyXMiniProblems.ProblemPhyXMini0117.RadioAntenna}
  \uses{def:physics:phyx-mini-0117:phyxminiproblems-problemphyxmini0117-frequencyquantity}
  One monochromatic radio antenna at the common reference time
\end{definition}
\begin{proof}
  This is the declared model definition of Radio Antenna.
\end{proof}
\begin{definition}[Two Antenna Interference Setup]
  \label{def:physics:phyx-mini-0117:phyxminiproblems-problemphyxmini0117-twoantennainterferencesetup}
  \lean{PhyXMiniProblems.ProblemPhyXMini0117.TwoAntennaInterferenceSetup}
  \uses{def:physics:phyx-mini-0117:phyxminiproblems-problemphyxmini0117-lengthquantity, def:physics:phyx-mini-0117:phyxminiproblems-problemphyxmini0117-frequencyquantity, def:physics:phyx-mini-0117:phyxminiproblems-problemphyxmini0117-speedquantity, def:physics:phyx-mini-0117:phyxminiproblems-problemphyxmini0117-figurepoint, def:physics:phyx-mini-0117:phyxminiproblems-problemphyxmini0117-antennalabel, def:physics:phyx-mini-0117:phyxminiproblems-problemphyxmini0117-answerchoice, def:physics:phyx-mini-0117:phyxminiproblems-problemphyxmini0117-radioantenna}
  The physical apparatus and its figure-derived quantities. The receiver position and both ray lengths are parameterized by the physical distance travelled from B. destructiveInterferenceAt records the observable to be predicted; it is not defined to make any displayed answer true
\end{definition}
\begin{proof}
  This is the declared model definition of Two Antenna Interference Setup.
\end{proof}
\begin{definition}[Radiates In Phase At Common Frequency]
  \label{def:physics:phyx-mini-0117:phyxminiproblems-problemphyxmini0117-radiatesinphaseatcommonfrequency}
  \lean{PhyXMiniProblems.ProblemPhyXMini0117.RadiatesInPhaseAtCommonFrequency}
  \uses{def:physics:phyx-mini-0117:phyxminiproblems-problemphyxmini0117-twoantennainterferencesetup}
  The two antennas emit in phase and at the common stated frequency
\end{definition}
\begin{proof}
  This is the declared model definition of Radiates In Phase At Common Frequency.
\end{proof}
\begin{definition}[Has Depicted Geometry]
  \label{def:physics:phyx-mini-0117:phyxminiproblems-problemphyxmini0117-hasdepictedgeometry}
  \lean{PhyXMiniProblems.ProblemPhyXMini0117.HasDepictedGeometry}
  \uses{def:physics:phyx-mini-0117:phyxminiproblems-problemphyxmini0117-lengthquantity, def:physics:phyx-mini-0117:phyxminiproblems-problemphyxmini0117-metersvalue, def:physics:phyx-mini-0117:phyxminiproblems-problemphyxmini0117-twoantennainterferencesetup}
  The categorical and coordinate layout read from the primary image. The coordinate chart is measured in metres: B is the origin, A lies on the positive vertical axis, and C lies on the positive horizontal axis. A receiver at nonnegative physical distance x from B has coordinates (x, 0). Thus the receiver ray BC is perpendicular to AB
\end{definition}
\begin{proof}
  This is the declared model definition of Has Depicted Geometry.
\end{proof}
\begin{definition}[Has Stated Readouts]
  \label{def:physics:phyx-mini-0117:phyxminiproblems-problemphyxmini0117-hasstatedreadouts}
  \lean{PhyXMiniProblems.ProblemPhyXMini0117.HasStatedReadouts}
  \uses{def:physics:phyx-mini-0117:phyxminiproblems-problemphyxmini0117-metersvalue, def:physics:phyx-mini-0117:phyxminiproblems-problemphyxmini0117-megahertzvalue, def:physics:phyx-mini-0117:phyxminiproblems-problemphyxmini0117-answerchoice, def:physics:phyx-mini-0117:phyxminiproblems-problemphyxmini0117-displayeddistancemeters, def:physics:phyx-mini-0117:phyxminiproblems-problemphyxmini0117-twoantennainterferencesetup}
  The numerical text and answer-table readouts. The wave frequency is 5.80 MHz, the antenna separation is 200 m, and radio propagation is modeled at Physlib's dimensionful speed of light. No interference behavior at any displayed distance occurs in this predicate
\end{definition}
\begin{proof}
  This is the declared model definition of Has Stated Readouts.
\end{proof}
\begin{definition}[Has Physical Parameters]
  \label{def:physics:phyx-mini-0117:phyxminiproblems-problemphyxmini0117-hasphysicalparameters}
  \lean{PhyXMiniProblems.ProblemPhyXMini0117.HasPhysicalParameters}
  \uses{def:physics:phyx-mini-0117:phyxminiproblems-problemphyxmini0117-lengthquantity, def:physics:phyx-mini-0117:phyxminiproblems-problemphyxmini0117-metersvalue, def:physics:phyx-mini-0117:phyxminiproblems-problemphyxmini0117-hertzvalue, def:physics:phyx-mini-0117:phyxminiproblems-problemphyxmini0117-meterspersecondvalue, def:physics:phyx-mini-0117:phyxminiproblems-problemphyxmini0117-antennalabel, def:physics:phyx-mini-0117:phyxminiproblems-problemphyxmini0117-answerchoice, def:physics:phyx-mini-0117:phyxminiproblems-problemphyxmini0117-twoantennainterferencesetup}
  Positivity and nonnegativity conditions for the physical quantities
\end{definition}
\begin{proof}
  This is the declared model definition of Has Physical Parameters.
\end{proof}
\begin{definition}[path Difference Meters]
  \label{def:physics:phyx-mini-0117:phyxminiproblems-problemphyxmini0117-pathdifferencemeters}
  \lean{PhyXMiniProblems.ProblemPhyXMini0117.pathDifferenceMeters}
  \uses{def:physics:phyx-mini-0117:phyxminiproblems-problemphyxmini0117-lengthquantity, def:physics:phyx-mini-0117:phyxminiproblems-problemphyxmini0117-metersvalue, def:physics:phyx-mini-0117:phyxminiproblems-problemphyxmini0117-twoantennainterferencesetup}
  The absolute geometric ray-path difference, read in metres
\end{definition}
\begin{proof}
  This is the declared model definition of path Difference Meters.
\end{proof}
\begin{definition}[Satisfies Free Space Two Source Laws]
  \label{def:physics:phyx-mini-0117:phyxminiproblems-problemphyxmini0117-satisfiesfreespacetwosourcelaws}
  \lean{PhyXMiniProblems.ProblemPhyXMini0117.SatisfiesFreeSpaceTwoSourceLaws}
  \uses{def:physics:phyx-mini-0117:phyxminiproblems-problemphyxmini0117-lengthquantity, def:physics:phyx-mini-0117:phyxminiproblems-problemphyxmini0117-metersvalue, def:physics:phyx-mini-0117:phyxminiproblems-problemphyxmini0117-hertzvalue, def:physics:phyx-mini-0117:phyxminiproblems-problemphyxmini0117-meterspersecondvalue, def:physics:phyx-mini-0117:phyxminiproblems-problemphyxmini0117-antennalabel, def:physics:phyx-mini-0117:phyxminiproblems-problemphyxmini0117-twoantennainterferencesetup, def:physics:phyx-mini-0117:phyxminiproblems-problemphyxmini0117-radiatesinphaseatcommonfrequency, def:physics:phyx-mini-0117:phyxminiproblems-problemphyxmini0117-pathdifferencemeters}
  The governing physical laws, stated without any answer-choice value: the vacuum dispersion relation is c = f * wavelength in SI readouts; each ray length is the Euclidean source-to-receiver distance; for in-phase sources, a nonnegative receiver distance is destructive iff its path difference is an odd multiple of one half-wavelength. The last clause is uniform over every nonnegative physical distance and does not assert that any of the four displayed candidates is destructive
\end{definition}
\begin{proof}
  This is the declared model definition of Satisfies Free Space Two Source Laws.
\end{proof}
\begin{lemma}[wavelength meters eq speed div frequency]
  \label{lem:physics:phyx-mini-0117:phyxminiproblems-problemphyxmini0117-wavelength-meters-eq-speed-div-frequency}
  \lean{PhyXMiniProblems.ProblemPhyXMini0117.wavelength_meters_eq_speed_div_frequency}
  \uses{def:physics:phyx-mini-0117:phyxminiproblems-problemphyxmini0117-metersvalue, def:physics:phyx-mini-0117:phyxminiproblems-problemphyxmini0117-hertzvalue, def:physics:phyx-mini-0117:phyxminiproblems-problemphyxmini0117-meterspersecondvalue, def:physics:phyx-mini-0117:phyxminiproblems-problemphyxmini0117-twoantennainterferencesetup, def:physics:phyx-mini-0117:phyxminiproblems-problemphyxmini0117-hasphysicalparameters, def:physics:phyx-mini-0117:phyxminiproblems-problemphyxmini0117-satisfiesfreespacetwosourcelaws}
  The vacuum dispersion law determines the wavelength from the stated frequency. This is a derived physical intermediate, not the requested receiver distance
\end{lemma}
\begin{proof}
  The conclusion follows from the governing relations and figure readouts named in the statement of wavelength meters eq speed div frequency.
\end{proof}
\begin{lemma}[path Difference Meters eq]
  \label{lem:physics:phyx-mini-0117:phyxminiproblems-problemphyxmini0117-pathdifferencemeters-eq}
  \lean{PhyXMiniProblems.ProblemPhyXMini0117.pathDifferenceMeters_eq}
  \uses{def:physics:phyx-mini-0117:phyxminiproblems-problemphyxmini0117-lengthquantity, def:physics:phyx-mini-0117:phyxminiproblems-problemphyxmini0117-metersvalue, def:physics:phyx-mini-0117:phyxminiproblems-problemphyxmini0117-twoantennainterferencesetup, def:physics:phyx-mini-0117:phyxminiproblems-problemphyxmini0117-hasdepictedgeometry, def:physics:phyx-mini-0117:phyxminiproblems-problemphyxmini0117-hasphysicalparameters, def:physics:phyx-mini-0117:phyxminiproblems-problemphyxmini0117-pathdifferencemeters, def:physics:phyx-mini-0117:phyxminiproblems-problemphyxmini0117-satisfiesfreespacetwosourcelaws}
  For a receiver x metres from B, the upper path has length sqrt(200\textasciicircum{}2 + x\textasciicircum{}2) while the lower path has length x. The statement is kept in terms of the physical separation so it follows from the depicted geometry before the 200 m readout is substituted
\end{lemma}
\begin{proof}
  The conclusion follows from the governing relations and figure readouts named in the statement of path Difference Meters eq.
\end{proof}
% --- Archon declaration topology end ---
% --- Archon physics formalization source end ---
